% ------------------------------------------------------------------------
%
\begin{frame}{Stacks in a functional setting}

  We show how to express linear data structures in a purely functional
  language and how to compute with them. The simplest of those
  structures is the \textbf{stack}.

  \bigskip

  A stack can be thought of as a finite series of items that can only
  be accessed sequentially from one end, called the \textbf{top},
  following the analogy with a stack of material objects.

\end{frame}

% ------------------------------------------------------------------------
%
\begin{frame}{Stacks in a functional setting}

  To define a stack in a functional language, one uses calls to
  undefined functions to model all the possible shapes a stack can
  take.

  \bigskip

  Here, a stack can be either empty or not.

  \bigskip

  Let \(\fun{Nil}()\) be the irreducible call that denotes an empty
  stack (that is, the function \fun{Nil} is undefined and its name is
  in uppercase to remind us of that fact).

  \bigskip

  If not empty, the stack then must contain a topmost (first) item,
  and a \textbf{substack}, that is, the stack made of all the items
  but the first.

  \bigskip

  Let \(\fun{Cons}(x,s)\) denote the non\hyp{}empty stack with the
  item~\(x\) on top and the \textbf{immediate substack}~\(s\).

\end{frame}

% ------------------------------------------------------------------------
%
\begin{frame}{Inductive definition of stacks}

  The previous definition is intuitive, but not formal yet.  Data
  structures are usually defined \textbf{inductively}.
  \begin{itemize}

    \item Let~\(\mathcal{T}\) be the set of all possible terms and
      \(\mathcal{S} \subset \mathcal{T}\) be the set of all stacks.

    \item Formally, \(\mathcal{S}\)~can by defined by
      \textbf{induction} as the smallest set such that
      \begin{enumerate}

        \item $\fun{Nil}() \in \mathcal{S}$;

        \item if \(x \in \mathcal{T}\) and \(s \in \mathcal{S}\), then
          \(\fun{Cons}(x,s) \in \mathcal{S}\).
      \end{enumerate}

  \end{itemize}

   We take the smallest set because we do not want in~\(\mathcal{S}\)
   terms that were not constructed by the inductive rules (this is
   called the \textbf{smallest fixed point} of the inductive
   relation). Another way to write the above is with inference rules:
   \begin{mathpar}
     \inferrule{}{\fun{Nil}() \in \mathcal{S}}
     \and
     \inferrule{x \in \mathcal{T}\\ s \in \mathcal{S}}%
               {\fun{Cons}(x,s) \in \mathcal{S}}
   \end{mathpar}

\end{frame}

% ------------------------------------------------------------------------
%
\begin{frame}{Appending a stack}

  \begin{itemize}

    \item For all \(s, t \in \mathcal{S}\), let us consider the
      functional program
      \begin{align*}
        \fun{cat}(\fun{Nil}(),t) & \xrightarrow{\alpha} t\\
        \fun{cat}(\fun{Cons}(x,s),t) &\xrightarrow{\smash\beta}
        \fun{Cons}(x,\fun{cat}(s,t))
      \end{align*}

    \item This is a functional program suitable for
      call\hyp{}by\hyp{}value reduction strategy, because the function
      calls in the patterns (that is, left\hyp{}hand sides) are
      irreducible.

    \item The right\hyp{}hand sides are stacks, so a call to
      \fun{cat} is a stack.

    \item Rule~\(\alpha\) says that, if the first stack is empty, the
      value is the second.

    \item In rule~\(\beta\), the first stack is not empty, and its top
      item is the top of the final stack~(\(x\)), the immediate
      substack~(\(s\)) being used in a call to \fun{cat} with the
      second stack~(\(t\)) unchanged.

  \end{itemize}

\end{frame}

% ------------------------------------------------------------------------
%
\begin{frame}{Appending a stack}

  This function appends the second stack at the bottom of a
  \textbf{copy} of the first. (The function name \fun{cat} stands
  for `catenation', and is slightly better than \fun{app}).

  \bigskip

  We could also say that this function appends a copy of the first
  stack on top of the other. But this is not the same as pushing the
  first stack on top of the other, like
  \begin{equation*}
    \fun{cat}(s,t) \rightarrow \fun{Cons}(s,t) \qquad
    \text{Wrong!}
  \end{equation*}

  This is \textbf{not} appending: here, the stack~\(s\) becomes an
  item of the resulting stack, but we want only the \textbf{contents}
  of~\(s\) to be pushed on top of~\(t\), not their container.

  \bigskip

  Analogy: You may think of an empty stack as a empty cardboard box,
  open topside, and every item pushed in it is like a book that you
  put on the topmost book, or at the bottom of the box if this is the
  first book.

\end{frame}

% ------------------------------------------------------------------------
%
\begin{frame}{Appending a stack}

  Let us set the abbreviations
  \begin{itemize}

    \item \(\el := \fun{Nil}()\),

    \item \(\cons{x}{s} := \fun{Cons}(x,s)\),

  \end{itemize}
  after the convention of the programming language \OCaml.

  \bigskip

  We can further abbreviate
  \begin{itemize}

    \item \(\cons{x_1}{\cons{x_2}{\cons{\dots}{\cons{x_n}{s}}}} :=
      \fun{Cons}(x_1, \fun{Cons}(x_2, \dots \fun{Cons}(x_n,s)))\)

    \item \([x] := \cons{x}{\el}\)

    \item \([x_1; x_2, \dots; x_n] :=
          \cons{x_1}{\cons{x_2}{\cons{x_3}{\cons{\ldots}{\cons{x_n}{\el}}}}}\).

  \end{itemize}

  Our system now becomes a bit more legible:
  \begin{align*}
    \fun{cat}(        \el,t) &\xrightarrow{\alpha} t\\
    \fun{cat}(\cons{x}{s},t) &\xrightarrow{\smash{\beta}}
    \cons{x}{\fun{cat}(s,t)}
  \end{align*}

\end{frame}

% ------------------------------------------------------------------------
%
\begin{frame}{Pattern matching}

  \begin{itemize}

    \item Let us compute \(\fun{cat}([1;2],[3])\).

    \item We need to find the first rule that is matched by this call.

    \item Remember that the arguments must be values before you do
      that.

    \item Intuitively, this means that we compare the call with the
      patterns, that is, the left\hyp{}hand sides of the rules, in
      order.

    \item Here, rule~\(\alpha\) is not matched because \([1;2] \neq
      []\).

    \item Rule~\(\beta\) is matched because \([1;2] = 1::[2]\) and
      equals the sub\hyp{}pattern \(x::s\) if we assign \(x := 1\) and
      \(s := [2]\). Also, trivially, \(t := [3]\) works.

    \item After the rule~\(\beta\) is selected, the assignments are
      used to replace the variables in the right\hyp{}hand sides.

    \item The whole evaluation is
    \begin{equation*}
      \fun{cat}([1;2],[3])
      \xrightarrow{\beta}
      \cons{1}{\fun{cat}([2],[3])}
      \xrightarrow{\beta}
      \cons{1}{\cons{2}{\fun{cat}(\el,[3])}}
      \xrightarrow{\alpha}
                    [1;2;3].
    \end{equation*}
  \end{itemize}

\end{frame}
